\subsection{Alcance de la implementación}

La implementación se concentró en incorporar un backend de almacenamiento
comprimido para el vector de estado dentro de TMFQSfullstate, manteniendo el modelo
de simulación de estado completo y preservando igualdad exacta respecto a la referencia
no comprimida. El objetivo no fue rediseñar el simulador en su totalidad, sino intervenir
la capa en la que se representan, recuperan y actualizan las amplitudes durante la
ejecución.

Bajo este alcance, el desarrollo incluyó tres elementos principales. El primero fue una
interfaz común para abstraer el almacenamiento del vector de estado y desacoplar el
registro cuántico de una representación específica. El segundo fue un backend
comprimido basado en Blosc2, encargado de la persistencia de bloques, la
descompresión cuando el acceso lo requiere y la escritura de los bloques modificados. El
tercero fue la integración de una estrategia de acceso y actualización compatible con la
aplicación de compuertas sobre amplitudes distribuidas en uno o varios bloques.

La implementación buscó, en consecuencia, verificar que la solución propuesta podía
insertarse en un simulador cuántico real con control suficiente sobre el uso de memoria,
la sobrecarga temporal y la exactitud numérica. Los resultados de esa evaluación se
presentan posteriormente; en esta sección interesa describir la organización del
desarrollo y la forma en que sus componentes cooperan durante la simulación.
